\subsection{Building Mathematical Optimization Models}\label{sec:introduction:building-mathematical-optimization-models}

The \textbf{goal} of this section is to \hl{learn how to go from a verbal description of a decision problem to a mathematical model.} The recurring structure is:
\begin{equation*}
    \text{Decision variables}
    +
    \text{Objective function}
    +
    \text{Constraints}
    +
    \text{Variable domains}
\end{equation*}
This is the basic language of mathematical optimization. The following examples illustrate how to use it.

\highspace
\begin{flushleft}
    \textcolor{Green3}{\faIcon{cubes} \textbf{The basic ingredients of an optimization model}}
\end{flushleft}
Suppose a problem asks:
\begin{center}
    ``\emph{What should we decide so that some objective is optimized while respecting certain restrictions?}''
\end{center}
We build the model in four steps.
\begin{enumerate}
    \item \important{Decision variables}. These represent the \hl{decisions under our control.} For example:
        \begin{equation*}
            x_j=\text{number of units of product }j
        \end{equation*}
        Or:
        \begin{equation*}
            x_j=
            \begin{cases}
                1 & \text{if investment }j\text{ is selected}\\
                0 & \text{otherwise}
            \end{cases}
        \end{equation*}
        The first modeling question should almost always be: ``\emph{\textbf{what do we need to decide?}}''. The answer usually tells us what our variables should be.

    \item \important{Objective function}. The objective tells us \hl{what ``best'' means.} Typical forms are:
        \begin{equation*}
            \max \text{ profit}
        \end{equation*}
        \begin{equation*}
            \min \text{ cost}
        \end{equation*}
        \begin{equation*}
            \min \text{ time}
        \end{equation*}
        \begin{equation*}
            \max \text{ return}
        \end{equation*}
        Mathematically:
        \begin{equation*}
            \max f(x)
            \qquad\text{or}\qquad
            \min f(x)
        \end{equation*}
        The objective is therefore a \hl{numerical measure of the quality of a feasible decision.}

    \item \important{Constraints}. Constraints represent the \hl{requirements that a valid solution must satisfy.} Examples include:
        \begin{equation*}
            \text{resource usage}\leq \text{available resource}
        \end{equation*}
        \begin{equation*}
            \text{number of selected items}\leq 5
        \end{equation*}
        Or:
        \begin{equation*}
            \text{distance from forbidden point}\geq 100
        \end{equation*}
        Constraints \hl{determine which candidate solutions are \textbf{feasible}.}

    \item \important{Variable domains}. We must also specify what \hl{values the variables are allowed to take.} For example:
        \begin{itemize}
            \item Continuous: $x_j\in\mathbb R$ or $x_j\geq 0$.
            \item Integer: $x_j\in\mathbb Z$.
            \item Binary: $x_j\in\{0,1\}$.
        \end{itemize}
        This is not a minor technical detail. The domain often completely changes the nature and difficulty of the optimization problem.
\end{enumerate}
So a model can be summarized schematically as:
\begin{equation*}
    \begin{aligned}
        \max/\min\quad & f(x)\\
        \text{s.t.}\quad & \text{constraints}\\
        & x\in\text{appropriate domain}
    \end{aligned}
\end{equation*}
Now let us see this structure in actual examples.

\highspace
\begin{examplebox}[: Production planning (continuous decision variables)]
    A company produces three electronic devices:
    \begin{equation*}
        D_1,\quad D_2,\quad D_3
    \end{equation*}
    Every product goes through three phases:
    \begin{enumerate}
        \item Assembly;
        \item Refinement;
        \item Quality control.
    \end{enumerate}
    The required processing times are:
    \begin{center}
        \begin{tabular}{@{} c | c c c @{}}
            \toprule
            & $D_1$ & $D_2$ & $D_3$\\
            \midrule
            Assembly & $80$ & $70$ & $120$\\
            Refinement & $70$ & $90$ & $20$\\
            Quality control & $40$ & $30$ & $20$\\
            \bottomrule
        \end{tabular}
    \end{center}
    All values are in minutes per unit. The available capacities are:
    \begin{center}
        \begin{tabular}{@{} c | c @{}}
            \toprule
            Phase & Available minutes\\
            \midrule
            Assembly & $30'000$\\
            Refinement & $25'000$\\
            Quality control & $18'000$\\
            \bottomrule
        \end{tabular}
    \end{center}
    The unit profits are:
    \begin{center}
        \begin{tabular}{@{} c | c c c @{}}
            \toprule
            & $D_1$ & $D_2$ & $D_3$\\
            \midrule
            Profit (k\euro) & $1.6$ & $1$ & $2$\\
            \bottomrule
        \end{tabular}
    \end{center}
    We assume everything produced can be sold. The goal is: \hl{choose production quantities to maximize total profit.}

    \highspace
    \textcolor{Green3}{\faIcon{sliders-h} \textbf{Step 1: Decision variables.}} \emph{What must the company decide?} How many units of each device to produce. Therefore define:
    \begin{equation*}
        x_1=\text{quantity of }D_1
    \end{equation*}
    \begin{equation*}
        x_2=\text{quantity of }D_2
    \end{equation*}
    \begin{equation*}
        x_3=\text{quantity of }D_3
    \end{equation*}
    Compactly:
    \begin{equation*}
        x_j=\text{number of units of }D_j
    \end{equation*}

    \highspace
    \textcolor{Green3}{\faIcon{bullseye} \textbf{Step 2: Objective function.}} Each unit of $D_1$ gives $1.6$ k\euro. So producing $x_1$ units gives:
    \begin{equation*}
        1.6x_1
    \end{equation*}
    Likewise:
    \begin{equation*}
        D_2:\quad 1x_2 = x_2
    \end{equation*}
    \begin{equation*}
        D_3:\quad 2x_3
    \end{equation*}
    Therefore \textbf{total profit} is:
    \begin{equation*}
        1.6x_1+x_2+2x_3
    \end{equation*}
    We want to maximize it:
    \begin{equation*}
        \boxed{
            \max z=1.6x_1+x_2+2x_3
        }
    \end{equation*}
    Notice the modeling pattern:
    \begin{equation*}
        \text{unit contribution}\times\text{quantity}
    \end{equation*}
    This appears constantly in OR.

    \highspace
    \textcolor{Green3}{\faIcon{tools} \textbf{Step 3: Assembly capacity.}} One unit of $D_1$ needs $80$ assembly minutes. Therefore $x_1$ units need:
    \begin{equation*}
        80x_1
    \end{equation*}
    Similarly:
    \begin{equation*}
        70x_2
    \end{equation*}
    For $D_2$, and:
    \begin{equation*}
        120x_3
    \end{equation*}
    For $D_3$. Hence \textbf{total assembly usage} is:
    \begin{equation*}
        80x_1+70x_2+120x_3
    \end{equation*}
    But only $30'000$ minutes are available:
    \begin{equation*}
        \boxed{
            80x_1+70x_2+120x_3\leq30000
        }
    \end{equation*}
    The logic is:
    \begin{equation*}
        \text{resource consumption}\leq\text{resource availability}
    \end{equation*}
    This is one of the \hl{most common constraint templates in optimization.}

    \highspace
    \textcolor{Green3}{\faIcon{cogs} \textbf{Step 4: Refinement capacity.}} Exactly the same reasoning gives:
    \begin{equation*}
        \boxed{
            70x_1+90x_2+20x_3\leq25000
        }
    \end{equation*}

    \highspace
    \textcolor{Green3}{\faIcon{check-circle} \textbf{Step 5: Quality-control capacity.}} Likewise:
    \begin{equation*}
        \boxed{
            40x_1+30x_2+20x_3\leq18000
        }
    \end{equation*}

    \highspace
    \textcolor{Green3}{\faIcon{shapes} \textbf{Step 6: Variable domains.}} Obviously, we cannot produce a negative quantity:
    \begin{equation*}
        x_1,x_2,x_3\geq0
    \end{equation*}
    The model used here allows fractional values:
    \begin{equation*}
        x_j\in\mathbb R_+
    \end{equation*}
    Where $\mathbb{R}_+$ denotes the set of non-negative real numbers (positive and zero). So, for example $x_1 = 125.7$ units of $D_1$ is mathematically allowed.

    \textcolor{Red2}{\faIcon{exclamation-triangle}} For actual indivisible electronic devices this may look strange, but this example is intentionally treated as a \textbf{continuous model} to introduce \emph{Linear Programming}. Later, integer restrictions could instead be imposed if needed.

    \highspace
    \textcolor{Green3}{\faIcon{sitemap} \textbf{Step 7: Complete model.}} Putting everything together, the complete mathematical optimization model is:
    \begin{equation*}
        \boxed{
            \begin{aligned}
                \max\quad
                &1.6x_1+x_2+2x_3\\[1mm]
                \text{s.t.}\quad
                &80x_1+70x_2+120x_3\leq30000\\
                &70x_1+90x_2+20x_3\leq25000\\
                &40x_1+30x_2+20x_3\leq18000\\
                &x_1,x_2,x_3\geq0.
            \end{aligned}
        }
    \end{equation*}
    Everything here is \textbf{linear}:
    \begin{itemize}
        \item Objective function: linear;
        \item Constraints: linear;
        \item Variables: continuous.
    \end{itemize}
    Therefore this is a \hl{Linear Programming problem (LP)}. We will formally define LP later.

    \highspace
    \begin{deepeningbox}[: How to read a constraint]
        We should develop the ability to look at:
        \begin{equation*}
            80x_1+70x_2+120x_3\leq30000
        \end{equation*}
        And immediately interpret it verbally: ``\emph{the total assembly time consumed by the production plan cannot exceed the available assembly time}''. Likewise:
        \begin{equation*}
            70x_1+90x_2+20x_3\leq25000
        \end{equation*}
        Means: ``\emph{the total refinement time cannot exceed the refinement capacity}''. This translation between mathematics and natural language is essential for modeling questions in the \textbf{exam}.
    \end{deepeningbox}
\end{examplebox}

\highspace
\begin{examplebox}[: Portfolio selection (binary decision variables)]\phantomsection\label{example:portfolio-selection}
    The second example is structurally different. An insurance company has eight possible investments:
    \begin{center}
        \begin{tabular}{@{} c c c @{}}
            \toprule
            \textbf{Investment} & \textbf{Sector/Type} & \textbf{Country} \\
            \midrule
            A & Automotive                & Germany \\[.3em]
            B & Automotive                & Italy   \\[.3em]
            C & ICT                       & USA     \\[.3em]
            D & ICT                       & Italy   \\[.3em]
            E & Real estate               & Italy   \\[.3em]
            F & Real estate               & France  \\[.3em]
            G & Short-term treasury bonds & Italy   \\[.3em]
            H & Long-term treasury bonds  & UK      \\
            \bottomrule
        \end{tabular}
    \end{center}
    Note that ICT stands for \emph{Information and Communication Technology}. It refers broadly to the technology sector: computing, software, telecommunications, electronics, IT services, etc. Instead, \emph{treasury bond} is a government-issued debt security. Very roughly, when we buy one, we are lending money to a government, and the government promises to repay us according to specified terms, generally with interest.

    \highspace
    Each investment has:
    \begin{itemize}
        \item A required capital $c_j$;
        \item An expected return $r_j$;
        \item A geographical location;
        \item A sector.
    \end{itemize}
    The available capital is:
    \begin{equation*}
        600\text{ k\euro}
    \end{equation*}
    There are also diversification rules:
    \begin{itemize}
        \item Select at most $5$ investments;
        \item Select at most $3$ investments in Italy;
        \item Select at most $3$ investments abroad.
    \end{itemize}
    The goal is: \hl{select investments maximizing expected return.}

    \highspace
    \textcolor{Green3}{\faIcon{question-circle} \textbf{Why the variable type is different.}} In production planning (previous example), asking:
    \begin{equation*}
        x_j=\text{how much of product }j\text{?}
    \end{equation*}
    Made sense. Here the decision is: ``\emph{do we select investment $j$ or not?}''. There are \textbf{only two possibilities}. So define:
    \begin{equation*}
        \boxed{
            x_j=
            \begin{cases}
                1 & \text{if investment }j\text{ is selected},\\
                0 & \text{otherwise}.
            \end{cases}
        }
    \end{equation*}
    Hence:
    \begin{equation*}
        x_j\in\{0,1\}
    \end{equation*}
    These are \textbf{binary variables}. This modeling trick is fundamental. Whenever the decision is of the form:
    \begin{equation*}
        \text{yes/no}
    \end{equation*}
    \begin{equation*}
        \text{open/closed}
    \end{equation*}
    \begin{equation*}
        \text{selected/not selected}
    \end{equation*}
    \begin{equation*}
        \text{build/do not build}
    \end{equation*}
    Binary variables are a \textbf{natural representation}.

    \highspace
    \textcolor{Green3}{\faIcon{bullseye} \textbf{Objective function.}} Investment $j$ requires capital $c_j$ and provides expected return rate $r_j$. Its expected monetary return is:
    \begin{equation*}
        c_jr_j
    \end{equation*}
    If it is \textbf{selected}:
    \begin{equation*}
        x_j=1
    \end{equation*}
    So its contribution is:
    \begin{equation*}
        c_jr_j
    \end{equation*}
    If it is \textbf{not selected}:
    \begin{equation*}
        x_j=0
    \end{equation*}
    So:
    \begin{equation*}
        c_jr_jx_j=0
    \end{equation*}
    Therefore \textbf{total expected return} is:
    \begin{equation*}
        \sum_{j=1}^{8}c_jr_jx_j
    \end{equation*}
    Where $8$ is the number of investments. And the objective becomes:
    \begin{equation*}
        \boxed{
            \max \sum_{j=1}^{8}c_jr_jx_j
        }
    \end{equation*}
    This illustrates another extremely \hl{common modeling device:}
    \begin{equation*}
        \text{coefficient}\times\text{binary variable}
    \end{equation*}
    Automatically activates or deactivates a contribution.

    \highspace
    \textcolor{Green3}{\faIcon{wallet} \textbf{Capital constraint.}} If investment $j$ is selected, it consumes $c_j$ units of capital. Therefore total invested capital is:
    \begin{equation*}
        \sum_{j=1}^{8}c_jx_j
    \end{equation*}
    Available capital is $600$ k\euro, so:
    \begin{equation*}
        \boxed{
            \sum_{j=1}^{8}c_jx_j\leq600
        }
    \end{equation*}

    \highspace
    \textcolor{Green3}{\faIcon{list-ol} \textbf{Maximum number of investments.}} Because $x_j=1$ exactly when investment $j$ is selected:
    \begin{equation*}
        \sum_{j=1}^{8}x_j
    \end{equation*}
    Counts the number of selected investments. Therefore:
    \begin{equation*}
        \boxed{
            \sum_{j=1}^{8}x_j\leq5
        }
    \end{equation*}
    Means: \emph{select at most five investments}. This is an incredibly \hl{useful property of binary variables:}
    \begin{equation}
        \sum_jx_j=\text{number of selected items}
    \end{equation}
    When every $x_j\in\{0,1\}$.

    \highspace
    \textcolor{Green3}{\faIcon{globe-europe} \textbf{Geographic diversification.}} The \textbf{investments} located in \textbf{Italy} are:
    \begin{equation*}
        B,D,E,G
    \end{equation*}
    Their variables are:
    \begin{equation*}
        x_2,x_4,x_5,x_7
    \end{equation*}
    At most three may be selected:
    \begin{equation*}
        \boxed{
            x_2+x_4+x_5+x_7\leq3
        }
    \end{equation*}
    Similarly, the \textbf{foreign investments} are:
    \begin{equation*}
        A,C,F,H
    \end{equation*}
    So:
    \begin{equation*}
        \boxed{
            x_1+x_3+x_6+x_8\leq3
        }
    \end{equation*}
    Again, the \hl{sum simply \textbf{counts selected investments belonging to a category}.} This is a basic modeling pattern worth memorizing:
    \begin{equation}
        \sum_{j\in S}x_j\leq k
    \end{equation}
    Means: ``\emph{select at most $k$ elements from set $S$}''. Likewise:
    \begin{equation*}
        \sum_{j\in S}x_j\geq k
    \end{equation*}
    Means: ``\emph{select at least $k$ elements from set $S$}''. And:
    \begin{equation*}
        \sum_{j\in S}x_j=k
    \end{equation*}
    Means: ``\emph{select exactly $k$ elements from set $S$}''.

    \highspace
    \textcolor{Green3}{\faIcon{sitemap} \textbf{Complete portfolio model.}} Putting everything together, the complete mathematical optimization model is:
    \begin{equation*}
        \boxed{
            \begin{aligned}
                \max\quad
                &\sum_{j=1}^{8}c_jr_jx_j\\[1mm]
                \text{s.t.}\quad
                &\sum_{j=1}^{8}c_jx_j\leq600\\
                &\sum_{j=1}^{8}x_j\leq5\\
                &x_2+x_4+x_5+x_7\leq3\\
                &x_1+x_3+x_6+x_8\leq3\\
                &x_j\in\{0,1\},
                \qquad j=1,\ldots,8
            \end{aligned}
        }
    \end{equation*}
    The \textbf{objective} and \textbf{constraints} are still \textbf{linear}. But the \textbf{variables} are \textbf{not continuous}:
    \begin{equation*}
        x_j\in\{0,1\}
    \end{equation*}
    Therefore this is an: Integer Linear Programming problem (ILP), or more precisely, a \textbf{binary linear optimization problem}.
\end{examplebox}

\begin{flushleft}
    \textcolor{Green3}{\faIcon{not-equal} \textbf{LP vs ILP: first important distinction}}
\end{flushleft}
The two examples are intentionally placed next to each other.
\begin{itemize}
    \item \textbf{Production planning}
        \begin{equation*}
            x_j\geq0,\qquad x_j\in\mathbb R
        \end{equation*}
        The variables can vary continuously. Therefore it is a \textbf{Linear Programming problem (LP)}.

    \item \textbf{Portfolio selection}
        \begin{equation*}
            x_j\in\{0,1\}
        \end{equation*}
        The variables are discrete. Therefore it is an \textbf{Integer Linear Programming problem (ILP)}.
\end{itemize}
Even though both models use linear equations and inequalities, the \textbf{variable domains are different}, and that \hl{difference matters enormously algorithmically.} We will later see that integer optimization is generally much harder than ordinary linear programming.

\begin{examplebox}[: Modeling logical implications with constraints]
    We continue with the \textbf{portfolio selection problem} (\autopageref{example:portfolio-selection}), but now add a logical rule: \hl{if at least one ICT investment is selected, then at least one treasury bond must also be selected.}

    \highspace
    From the table on \autopageref{example:portfolio-selection}, we know that the ICT investments are:
    \begin{equation*}
        C,D
    \end{equation*}
    With variables:
    \begin{equation*}
        x_3,x_4
    \end{equation*}
    The treasury bonds are:
    \begin{equation*}
        G,H
    \end{equation*}
    With variables:
    \begin{equation*}
        x_7,x_8
    \end{equation*}
    Recall that all variables are binary:
    \begin{equation*}
        x_j\in\{0,1\}
    \end{equation*}
    So we want to \textbf{translate the logical statement}:
    \begin{equation*}
        (x_3=1\ \text{or}\ x_4=1)
        \Longrightarrow
        (x_7=1\ \text{or}\ x_8=1)
    \end{equation*}
    \textbf{Into a linear constraint}.

    \highspace
    \textcolor{Green3}{\faIcon{list-ol} \textbf{Counting selected investments.}} Because the variables are binary:
    \begin{equation*}
        x_3+x_4
    \end{equation*}
    \hl{Counts the number of selected ICT investments.} Therefore:
    \begin{equation*}
        x_3+x_4=
        \begin{cases}
            0 & \text{no ICT investment}\\
            1 & \text{one ICT investment}\\
            2 & \text{both ICT investments}
        \end{cases}
    \end{equation*}
    Similarly:
    \begin{equation*}
        x_7+x_8
    \end{equation*}
    \hl{Counts selected treasury bonds.} We need:
    \begin{equation*}
        x_3+x_4>0
        \quad\Longrightarrow\quad
        x_7+x_8\geq1.
    \end{equation*}
    The model uses:
    \begin{equation*}
        \boxed{
            \frac{x_3+x_4}{2}\leq x_7+x_8
        }
    \end{equation*}
    Let's understand \textbf{why this works}.

    \highspace
    \textcolor{Green3}{\faIcon{times-circle} \textbf{No ICT investment.}} Suppose:
    \begin{equation*}
        x_3=x_4=0
    \end{equation*}
    Then:
    \begin{equation*}
        \frac{x_3+x_4}{2}=0
    \end{equation*}
    The constraint becomes:
    \begin{equation*}
        0\leq x_7+x_8
    \end{equation*}
    This is always true. Therefore, if we select no ICT investment, there is \textbf{no obligation} to select a treasury bond. Exactly what we want.

    \highspace
    \textcolor{Green3}{\faIcon{check-circle} \textbf{One ICT investment.}} Suppose:
    \begin{equation*}
        x_3+x_4=1
    \end{equation*}
    Then:
    \begin{equation*}
        \frac{x_3+x_4}{2}=\frac{1}{2}
    \end{equation*}
    The constraint becomes:
    \begin{equation*}
        \frac{1}{2} \leq x_7+x_8
    \end{equation*}
    But $x_7+x_8$ is integer because $x_7,x_8$ are binary. Its possible values are:
    \begin{equation*}
        0,1,2
    \end{equation*}
    Therefore:
    \begin{equation*}
        x_7+x_8 \geq \frac{1}{2}
    \end{equation*}
    Actually forces:
    \begin{equation*}
        x_7+x_8\geq1
    \end{equation*}
    Because $x_7+x_8$ cannot be $0$ and the next possible value is $1$. So \textbf{at least one treasury bond must be selected}.

    \highspace
    \textcolor{Green3}{\faIcon{check-double} \textbf{Both ICT investments.}} Suppose:
    \begin{equation*}
        x_3=x_4=1
    \end{equation*}
    Then:
    \begin{equation*}
        \frac{x_3+x_4}{2}=1
    \end{equation*}
    Therefore:
    \begin{equation*}
        1\leq x_7+x_8
    \end{equation*}
    Again \textbf{at least one treasury bond} must be selected. Notice that selecting both ICT investments does \textbf{not} require two treasury bonds. The logical rule only says: ``\emph{if \textbf{any} ICT investment is present, select \textbf{at least once} bond}''. So the formulation is correct.

    \highspace
    \textcolor{Green3}{\faIcon{exchange-alt} \textbf{Equivalent form.}} We can multiply both sides by $2$:
    \begin{equation*}
        x_3+x_4\leq2(x_7+x_8)
    \end{equation*}
    Which may make the logic easier to see.
\end{examplebox}

\begin{flushleft}
    \textcolor{Green3}{\faIcon{project-diagram} \textbf{A general modeling pattern for \emph{Logical Implications}}}
\end{flushleft}
From the previous example, we can extract a \textbf{general modeling pattern}. Suppose $S$ is a group containing $m$ binary variables, and we want: \hl{if at least one element of $S$ is selected, then at least one element of $T$ must be selected.} A formulation is:
\begin{equation}
    \frac{1}{m}\sum_{j\in S}x_j
    \leq
    \sum_{k\in T}x_k
\end{equation}
Or equivalently:
\begin{equation}
    \sum_{j\in S}x_j
    \leq
    m\sum_{k\in T}x_k
\end{equation}
This kind of construction will become very important in \textbf{Integer Linear Programming (ILP)}, because binary variables allow us to encode logical decisions using algebra.

\begin{examplebox}[: Facility location (nonlinear models)]
    The third example is geometrically very different. There are three oil pits:
    \begin{equation*}
        A=(0,0)
    \end{equation*}
    \begin{equation*}
        B=(300,0)
    \end{equation*}
    \begin{equation*}
        C=(240,300)
    \end{equation*}
    We must choose \textbf{where to build a refinery}, which will be connected to the pits by pipelines (i.e., straight lines). There is also a point:
    \begin{equation*}
        D=(100,200)
    \end{equation*}
    And the refinery must be \hl{at least $100$ km away from $D$.}

    \highspace
    The oil pits will be connected directly to the refinery by pipelines. The important detail is that the pipeline cost is proportional to the \textbf{square of its length}. We want to \hl{determine the refinery location minimizing total pipeline cost.}

    \highspace
    \begin{deepeningbox}[: What does ``\emph{cost proportional to the square of the length}'' mean?]
        Suppose a pipeline has length $L$. The exercise assumes its cost behaves like:
        \begin{equation*}
            \text{cost}=kL^2
        \end{equation*}
        For some positive constant $k$. So if one pipeline has length (i.e., distance from one oil pit to the refinery):
        \begin{equation*}
            L=10
        \end{equation*}
        Its cost is proportional to:
        \begin{equation*}
            10^2=100
        \end{equation*}
        If another has length:
        \begin{equation*}
            L=20
        \end{equation*}
        Its cost is proportional to:
        \begin{equation*}
            20^2=400
        \end{equation*}
        So doubling the length does not merely double the cost:
        \begin{equation*}
            20=2\cdot10
        \end{equation*}
        But:
        \begin{equation*}
            20^2=4\cdot10^2
        \end{equation*}
        Hence the cost becomes four times as large. The exercise gives this as a modeling assumption. We do not need to justify physically why pipeline construction behaves exactly like $L^2$; we simply model what the problem states.
    \end{deepeningbox}

    \highspace
    \textcolor{Green3}{\faIcon{crosshairs} \textbf{Step 1: Decision variables.}} \emph{What do we decide?} The refinery coordinates. So we define two variables:
    \begin{equation*}
        x_1=\text{x-coordinate of the refinery}
    \end{equation*}
    \begin{equation*}
        x_2=\text{y-coordinate of the refinery}
    \end{equation*}
    Therefore:
    \begin{equation*}
        (x_1,x_2)\in\mathbb R^2
    \end{equation*}
    Unlike the portfolio problem, these are \textbf{continuous variables} (i.e., they can take any real value).

    \highspace
    \textcolor{Green3}{\faIcon{ruler} \textbf{Step 2: Distance to each oil pit.}} Recall the Euclidean distance between points:
    \begin{equation*}
        (x_1,x_2)
    \end{equation*}
    And:
    \begin{equation*}
        (a,b)
    \end{equation*}
    Is:
    \begin{equation*}
        d=\sqrt{(x_1-a)^2+(x_2-b)^2}
    \end{equation*}
    But pipeline cost is proportional to the \textbf{square} of the distance. Therefore:
    \begin{equation*}
        d^2 = \left(\sqrt{(x_1-a)^2+(x_2-b)^2}\right)^{2} = (x_1-a)^2+(x_2-b)^2
    \end{equation*}
    This is convenient because the square root disappears.
    \begin{itemize}
        \item \textbf{Pipeline from $A$}. Since $A=(0,0)$, the squared length is:
            \begin{equation*}
                (x_1-0)^2+(x_2-0)^2
            \end{equation*}
            Thus:
            \begin{equation*}
                x_1^2+x_2^2
            \end{equation*}

        \item \textbf{Pipeline from $B$}. Since $B=(300,0)$, the squared length is:
            \begin{equation*}
                (x_1-300)^2+x_2^2
            \end{equation*}

        \item \textbf{Pipeline from $C$}. Since $C=(240,300)$, the squared length is:
            \begin{equation*}
                (x_1-240)^2+(x_2-300)^2
            \end{equation*}
    \end{itemize}

    \highspace
    \textcolor{Green3}{\faIcon{bullseye} \textbf{Step 3: Objective function.}} \hl{Total cost is proportional to the sum of these squared distances.} Therefore:
    \begin{equation*}
        \boxed{
            \begin{aligned}
                \min z={}&
                \left[x_1^2+x_2^2\right]\\
                &+\left[(x_1-300)^2+x_2^2\right]\\
                &+\left[(x_1-240)^2+(x_2-300)^2\right]
            \end{aligned}
        }
    \end{equation*}
    This already differs fundamentally from our previous models. For production planning, the objective looked like:
    \begin{equation*}
        1.6x_1+x_2+2x_3
    \end{equation*}
    It was \textbf{linear}. Here we have terms such as $x_1^2$ and $x_2^2$. Therefore the objective is \textbf{nonlinear}. More precisely, it is a \textbf{quadratic function}.

    \highspace
    \textcolor{Green3}{\faIcon{ban} \textbf{Step 4: Forbidden area around $D$.}} The refinery must remain at least $100$ km away from $D$:
    \begin{equation*}
        D=(100,200)
    \end{equation*}
    The distance between the refinery and $D$ is:
    \begin{equation*}
        \sqrt{(x_1-100)^2+(x_2-200)^2}
    \end{equation*}
    Therefore the requirement becomes:
    \begin{equation*}
        \boxed{
            \sqrt{(x_1-100)^2+(x_2-200)^2}\geq100
        }
    \end{equation*}
    To enhance geometric intuition, since both sides are nonnegative, we can square both sides:
    \begin{equation*}
        (x_1-100)^2+(x_2-200)^2\geq10000
    \end{equation*}
    This is a \textbf{nonlinear constraint} because it contains squared decision variables. Geometrically, the equation:
    \begin{equation*}
        \sqrt{(x_1-100)^2+(x_2-200)^2}=100
    \end{equation*}
    Describes a circle centered at $(100,200)$ with radius $100$. The inequality $\geq 100$ means the refinery must lie \textbf{\emph{on} or \emph{outside} that circle}. It cannot lie inside the forbidden disk.

    \highspace
    \textcolor{Green3}{\faIcon{sitemap} \textbf{Complete facility-location model.}} Putting everything together, we have:
    \begin{equation*}
        \boxed{
            \begin{aligned}
                \min\quad&
                x_1^2+x_2^2\\
                &+(x_1-300)^2+x_2^2\\
                &+(x_1-240)^2+(x_2-300)^2\\[1mm]
                \text{s.t.}\quad&
                (x_1-100)^2+(x_2-200)^2\geq10000\\
                &x_1,x_2\in\mathbb{R}
            \end{aligned}
        }
    \end{equation*}
    This is a Nonlinear Programming problem (NLP) because the objective and constraints are not all linear.

    \highspace
    \textcolor{Green3}{\faIcon{question-circle} \textbf{Why this is different from LP.}} A linear function can contain:
    \begin{equation*}
        3x_1+7x_2-2x_3
    \end{equation*}
    But it cannot contain: $x_1^2$, $x_1x_2$, $\sqrt{x_1}$, $e^{x_1}$, etc. So LP variables appear only linearly, whereas here:
    \begin{equation*}
        (x_1-300)^2
    \end{equation*}
    Contains squared decision variables. Hence it is a nonlinear model.
\end{examplebox}

\highspace
\begin{takeawaysbox}
    The most important modeling template is:
    \begin{equation*}
        \boxed{
            \begin{aligned}
                \text{Decision variables}&:\quad \text{What do I control?}\\
                \text{Objective}&:\quad \text{What do I optimize?}\\
                \text{Constraints}&:\quad \text{What must be respected?}\\
                \text{Domains}&:\quad \text{What values can the variables take?}
            \end{aligned}
        }
    \end{equation*}
    For \textbf{resource constraints}:
    \begin{equation*}
        \boxed{
            \text{total consumption}
            \leq
            \text{available capacity}
        }
    \end{equation*}
    For \textbf{binary variables}:
    \begin{equation*}
        x_j=
        \begin{cases}
            1 & \text{selected}\\
            0 & \text{not selected}
        \end{cases}
    \end{equation*}
    And therefore:
    \begin{equation*}
        \boxed{\sum_jx_j=\text{number of selected items}}
    \end{equation*}
    Finally:
    \begin{center}
        \begin{tabular}{@{} c c c @{}}
            \toprule
            & Production planning & Portfolio selection\\
            \midrule
            Variables & continuous & binary\\
            Objective & linear & linear\\
            Constraints & linear & linear\\
            Class & LP & ILP\\
            \bottomrule
        \end{tabular}
    \end{center}

    \highspace
    For \textbf{binary modeling}, remember that \hl{logical statements can often be encoded through inequalities.} In the previous examples:
    \begin{equation*}
        \text{ICT selected}
        \Rightarrow
        \text{treasury bond selected}
    \end{equation*}
    Mathematically, this was expressed as:
    \begin{equation*}
        x_3+x_4 \leq 2(x_7+x_8)
    \end{equation*}
    For geometric modeling, remember that the Euclidean distance is:
    \begin{equation*}
        d((x_1,x_2),(a,b))
        =
        \sqrt{(x_1-a)^2+(x_2-b)^2}
    \end{equation*}
    And squared-distance costs naturally produce \textbf{quadratic objectives}. Most importantly, this section has now introduced the four fundamental pieces we will repeatedly use:
    \begin{equation*}
        \text{variables}
        +
        \text{objective}
        +
        \text{constraints}
        +
        \text{domains}
    \end{equation*}
\end{takeawaysbox}